\Needspace{22\baselineskip}\SourceLeaf{31}
\[
\begin{array}{r|rrr}
n & a & b & c \\
\hline
2 & -1 & 1 & 0 \\
3 & 2 & 0 & -1 \\
4 & 1 & 1 & -1 \\
5 & 1 & -4 & 2 \\
6 & 1 & -6 & 3 \\
7 & 1 & -1 & 0 \\
8 & -2 & 0 & 1 \\
9 & -2 & 0 & 1 \\
10 & 1 & 6 & -3 \\
11 & 1 & 4 & -2 \\
12 & -107 & 33 & 6 \\
13 & 4 & 3 & -2 \\
14 & 1 & 2 & -1 \\
15 & 1 & -30 & 12 \\
16 & 1 & 50 & -20 \\
17 & 18 & -7 & 0 \\
19 & -8 & 3 & 0 \\
37 & 10 & -3 & 0 \\
\end{array}
\]

\Needspace{12\baselineskip}
\begingroup\small

\begin{center}\textbf{REMARKS ON THE FRAGMENTS ON THE THEORY OF ALGEBRAIC INTEGERS
FORMED FROM A CUBE ROOT.}\end{center}\par\noindent\hspace*{\parindent}The developments reproduced here date from about the same time as
the preceding »Zur Theorie der cubischen Reste«. According to a
diary entry, \textsc{Gauss} began the particular investigation of solving the
equation $x^3+ny^3+nnz^3-3nxyz=1$ on 23 Dec. 1808. This provides
further evidence that algebraic integers appeared at such an early
date.

In the language of more recent number theory, the subject here is the integers
of the real cubic field defined by the pure equation $x^3=n$.
Here $t$ denotes an integer of this field, $\varphi t$ the
norm of $t$, and $Ft$ the product, also in the field, of the numbers conjugate to
$t$, which we may call $t_1$, $t_2$.

\[
\varphi t = t t_1 t_2, \qquad Ft = t_1 t_2.
\]
